\documentclass{article}

\usepackage[T1]{fontenc}
\usepackage{geometry}
\usepackage{listings}
\usepackage{hyperref}

\setlength{\parindent}{0pt}
\setlength{\parskip}{0.6em}

\geometry{margin=1in}

\lstdefinelanguage{EBNF}{
  morekeywords={},
  sensitive=true,
  morecomment=[l]{//},
  morestring=[b]"
}

\lstset{
  basicstyle=\ttfamily\small,
  frame=single,
  columns=fullflexible,
  keepspaces=true
}

\title{Exchange Command Language Specification}
\author{Will Lucas}
\date{}

\begin{document}

\maketitle
\tableofcontents

\section{Overview}

This document specifies the textual command language accepted by the
exchange simulator when it is operating in \textit{manual mode}.

This allows the user to intervene and act as any participant. The user can:

\begin{itemize}
    \item add instruments to the exchange;
    \item add participants to the exchange;
    \item list the exchange's instruments and participants;
    \item submit buy and sell orders on behalf of participants;
    \item cancel existing orders on behalf of their owning participants;
    \item inspect a participant's available and reserved assets;
    \item inspect an instrument's order book;
    \item inspect the exchange's completed trades;
    \item inspect the current simulation time;
    \item control the speed of simulation time;
    \item and display a price graph for an instrument.
\end{itemize}

Commands which alter a participant's orders must explicitly identify the
participant using the form:

\begin{lstlisting}
AS <participant_id> <command>
\end{lstlisting}

For example:

\begin{lstlisting}
AS ALICE BUY AAPL LIMIT PRICE 10490 QUANTITY 7
\end{lstlisting}

Here, \texttt{ALICE} is the participant's unique
\texttt{participant\_id}.

Administrative commands and read-only commands do not require an
\texttt{AS} prefix.

\section{Lexical Structure}

Commands are case-sensitive. All alphabetic characters in keywords and
identifiers must be uppercase. For example, the following command is valid:

\begin{lstlisting}
SHOW BOOK AAPL
\end{lstlisting}

The following command is invalid:

\begin{lstlisting}
show book aapl
\end{lstlisting}

Tokens are separated by one or more whitespace characters.

\subsection{Numbers}

Prices, quantities, and order identifiers are represented by positive
integers.

The exchange stores prices using integer price units rather than
floating-point numbers. For example, if the configured price unit is one
cent, then:

\begin{lstlisting}
PRICE 10490
\end{lstlisting}

represents a price of \$104.90.

A quantity must be greater than zero.

\section{Grammar}

The command language is specified below using Extended Backus--Naur Form
(EBNF).

Terminal symbols are enclosed in double quotation marks. Braces denote
zero or more repetitions, and brackets group alternatives.

\begin{lstlisting}[language=EBNF]
command =
      add-instrument-command
    | add-participant-command
    | list-command
    | set-time-multiplier-command
    | fast-forward-command
    | participant-command
    | show-command
    ;

add-instrument-command =
    "ADD", "INSTR", symbol,
    "PRICE", price, "VOLUME", quantity
    ;

add-participant-command =
    "ADD", "PARTICIPANT", participant-id
    ;

list-command =
      "LIST", "INSTR"
    | "LIST", "PARTICIPANT"
    ;

set-time-multiplier-command =
    "SET", "TIME", "MULTIPLIER", nonnegative-integer
    ;

fast-forward-command =
    "FAST", "FORWARD", nonnegative-integer
    ;

participant-command =
    "AS", participant-id,
    (
          place-order-command
        | cancel-order-command
    )
    ;

place-order-command =
    side, symbol, order-specification,
    "QUANTITY", quantity
    ;

side =
      "BUY"
    | "SELL"
    ;

order-specification =
      "MARKET"
    | "LIMIT", "PRICE", price
    ;

cancel-order-command =
    "CANCEL", "ORDER", order-id
    ;

show-command =
      "SHOW", "BOOK", symbol
    | "SHOW", "PARTICIPANT", participant-id
    | "SHOW", "TRADES"
    | "SHOW", "TIME"
    | "SHOW", "GRAPH", symbol
    ;

participant-id =
    identifier
    ;

symbol =
    identifier
    ;

price =
    positive-integer
    ;

quantity =
    positive-integer
    ;

nonnegative-integer =
    digit, { digit }
    ;

order-id =
    positive-integer
    ;

identifier =
    letter, { letter | digit | "_" }
    ;

positive-integer =
    nonzero-digit, { digit }
    ;

letter =
      "A" | "B" | "C" | "D" | "E" | "F"
    | "G" | "H" | "I" | "J" | "K" | "L"
    | "M" | "N" | "O" | "P" | "Q" | "R"
    | "S" | "T" | "U" | "V" | "W" | "X"
    | "Y" | "Z"
    ;

digit =
      "0" | nonzero-digit
    ;

nonzero-digit =
      "1" | "2" | "3" | "4" | "5"
    | "6" | "7" | "8" | "9"
    ;
\end{lstlisting}

\section{Command Reference}

\subsection{\texttt{ADD INSTR}}

\subsubsection*{Syntax}

\begin{lstlisting}
ADD INSTR <symbol> PRICE <price> VOLUME <volume>
\end{lstlisting}

\subsubsection*{Description}

Adds a new tradable instrument to the exchange and issues its initial volume.

The supplied symbol must be a valid identifier and must not already
belong to another instrument. The price and volume must be positive integers.

The issued volume is divided as evenly as possible between the simulation's
configured initial position holders. In the default configuration these are
the market-maker and long-term-holder participants. The supplied price is
recorded as their acquisition cost. No orders are submitted automatically.

The supplied price also becomes the instrument's initial hidden fundamental
value. Its initial sentiment is neutral and its initial volatility is 0.1 per
cent. Hidden market state is maintained by the simulation and is not exposed
through the exchange's public order book.

\subsubsection*{Example}

\begin{lstlisting}
ADD INSTR AAPL PRICE 10000 VOLUME 100
\end{lstlisting}

After this command succeeds, 100 units of \texttt{AAPL} are allocated across
the configured initial position holders at an acquisition cost of 10000 ticks.

\subsubsection*{Possible errors}

The command will fail if:

\begin{itemize}
    \item the symbol is not a valid identifier;
    \item the symbol is a reserved keyword; or
    \item an instrument with the same symbol already exists;
    \item the price or volume is not positive.
\end{itemize}

\subsection{\texttt{ADD PARTICIPANT}}

\subsubsection*{Syntax}

\begin{lstlisting}
ADD PARTICIPANT <participant_id>
\end{lstlisting}

\subsubsection*{Description}

Registers a new participant with the exchange.

The participant identifier uniquely identifies the owner of any orders
submitted using the corresponding \texttt{AS} prefix.

\subsubsection*{Example}

\begin{lstlisting}
ADD PARTICIPANT ALICE
\end{lstlisting}

After this command succeeds, orders may be submitted on behalf of
\texttt{ALICE}.

\subsubsection*{Possible errors}

The command will fail if:

\begin{itemize}
    \item the participant identifier is invalid;
    \item the participant identifier is a reserved keyword; or
    \item a participant with the same identifier already exists.
\end{itemize}

\subsection{\texttt{LIST INSTR}}

\subsubsection*{Syntax}

\begin{lstlisting}
LIST INSTR
\end{lstlisting}

\subsubsection*{Description}

Displays the symbol and total issued supply of every instrument currently
registered with the exchange. Instruments are displayed in registration
order.

If no instruments exist, the command displays an empty list rather than
returning an error.

\subsubsection*{Example}

\begin{lstlisting}
LIST INSTR
\end{lstlisting}

\subsection{\texttt{LIST PARTICIPANT}}

\subsubsection*{Syntax}

\begin{lstlisting}
LIST PARTICIPANT
\end{lstlisting}

\subsubsection*{Description}

Displays the identifiers, available cash balances, and net worth of all participants
currently registered with the exchange. Participants are displayed in
registration order. Balances are formatted as dollar amounts; internally,
one integer cash unit represents one cent.

Net worth is the participant's total cash plus the current marked value of
all available and reserved positions.

If no participants exist, the command displays an empty list rather than
returning an error.

\subsubsection*{Example}

\begin{lstlisting}
LIST PARTICIPANT
\end{lstlisting}

\subsection{\texttt{AS \ldots BUY/SELL \ldots MARKET}}

\subsubsection*{Syntax}

\begin{lstlisting}
AS <participant_id> BUY  <symbol> MARKET QUANTITY <quantity>
AS <participant_id> SELL <symbol> MARKET QUANTITY <quantity>
\end{lstlisting}

\subsubsection*{Description}

Submits a market order on behalf of the specified participant.

A market order attempts to trade immediately against compatible resting
orders on the opposite side of the order book.

A buy market order matches against the lowest-priced sell orders. A sell
market order matches against the highest-priced buy orders.

Orders at the same price are matched according to time priority.

The remainders of unfilled market orders are cancelled immediately.

\subsubsection*{Examples}

\begin{lstlisting}
AS ALICE BUY AAPL MARKET QUANTITY 7
AS BOB SELL AAPL MARKET QUANTITY 3
\end{lstlisting}

\subsubsection*{Possible errors}

The command will fail if:

\begin{itemize}
    \item the participant does not exist;
    \item the instrument does not exist;
    \item the quantity is not a positive integer; or
    \item the participant is not permitted to submit the order.
\end{itemize}

\subsection{\texttt{AS \ldots BUY/SELL \ldots LIMIT}}

\subsubsection*{Syntax}

\begin{lstlisting}
AS <participant_id> BUY  <symbol> LIMIT PRICE <price> QUANTITY <quantity>
AS <participant_id> SELL <symbol> LIMIT PRICE <price> QUANTITY <quantity>
\end{lstlisting}

\subsubsection*{Description}

Submits a limit order on behalf of the specified participant.

A limit order may execute only at its specified price or at a better
price.

A buy limit order may match against sell orders whose prices are less
than or equal to its limit price. A sell limit order may match against
buy orders whose prices are greater than or equal to its limit price.

If the order is not completely filled immediately, its remaining
quantity rests on the order book.

Resting orders are prioritised first by price and then by submission
time.

Orders submitted through the command interface do not have an automatic
expiry and remain active until filled or cancelled. Automated strategies may
set an expiry timestamp directly on their order requests.

\subsubsection*{Examples}

\begin{lstlisting}
AS ALICE BUY AAPL LIMIT PRICE 10490 QUANTITY 7
AS BOB SELL AAPL LIMIT PRICE 10500 QUANTITY 4
\end{lstlisting}

\subsubsection*{Possible errors}

The command will fail if:

\begin{itemize}
    \item the participant does not exist;
    \item the instrument does not exist;
    \item the price is not a positive integer;
    \item the quantity is not a positive integer; or
    \item the participant is not permitted to submit the order.
\end{itemize}

\subsection{\texttt{AS \ldots CANCEL ORDER}}

\subsubsection*{Syntax}

\begin{lstlisting}
AS <participant_id> CANCEL ORDER <order_id>
\end{lstlisting}

\subsubsection*{Description}

Cancels the remaining unfilled quantity of an active order.

The participant identifier is mandatory. A participant may cancel only
an order which it owns.

Completed, fully filled, previously cancelled, or otherwise inactive
orders cannot be cancelled.

\subsubsection*{Example}

\begin{lstlisting}
AS ALICE CANCEL ORDER 42
\end{lstlisting}

This requests cancellation of order \texttt{42}, provided that the order
belongs to \texttt{ALICE} and remains active.

\subsubsection*{Possible errors}

The command will fail if:

\begin{itemize}
    \item the participant does not exist;
    \item the order does not exist;
    \item the order does not belong to the specified participant; or
    \item the order is no longer active.
\end{itemize}

\subsection{\texttt{SHOW BOOK}}

\subsubsection*{Syntax}

\begin{lstlisting}
SHOW BOOK <symbol>
\end{lstlisting}

\subsubsection*{Description}

Displays the current order book for an instrument.

The output shows active buy and sell orders, grouped or ordered by
price level. Buy orders are displayed from highest price
to lowest price, while sell orders are displayed from lowest price
to highest price.

The display may include:

\begin{itemize}
    \item the best bid;
    \item the best ask;
    \item the spread;
    \item aggregate quantity at each price level;
    \item individual order identifiers;
    \item participant identifiers; and
    \item order submission times.
\end{itemize}

\subsubsection*{Example}

\begin{lstlisting}
SHOW BOOK AAPL
\end{lstlisting}

\subsubsection*{Possible errors}

The command will fail if the specified instrument does not exist.

\subsection{\texttt{SHOW PARTICIPANT}}

\subsubsection*{Syntax}

\begin{lstlisting}
SHOW PARTICIPANT <participant_id>
\end{lstlisting}

\subsubsection*{Description}

Displays an accounting snapshot for the specified participant. The cash
summary contains available cash, cash reserved for active buy orders, their
combined total, unrealised gain or loss, and net worth. Positive unrealised
gains are displayed in green and negative unrealised gains in red.

Positions are marked using the exchange's current reference price. Unrealised
gain or loss is the difference between their marked value and average cost
basis. Net worth is total cash plus marked position value.

For each held instrument, the position table contains the quantity available
to sell, the quantity reserved for active sell orders, and their combined
total. It also displays the average acquisition cost per unit and the current
mark price used for valuation.

\subsubsection*{Example}

\begin{lstlisting}
SHOW PARTICIPANT ALICE
\end{lstlisting}

\subsubsection*{Possible errors}

The command fails if the specified participant does not exist.

\subsection{\texttt{SHOW TRADES}}

\subsubsection*{Syntax}

\begin{lstlisting}
SHOW TRADES
\end{lstlisting}

\subsubsection*{Description}

Displays trades completed during the current exchange session.

Trades are displayed in execution order.

Each displayed trade includes:

\begin{itemize}
    \item the trade identifier;
    \item the instrument symbol;
    \item the execution price;
    \item the executed quantity;
    \item the buy participant;
    \item the sell participant;
    \item the associated order identifiers; and
    \item the execution timestamp.
\end{itemize}

\subsubsection*{Example}

\begin{lstlisting}
SHOW TRADES
\end{lstlisting}

If no trades have occurred, the command displays an empty trade
history rather than returning an error.

\subsection{\texttt{SHOW TIME}}

\subsubsection*{Syntax}

\begin{lstlisting}
SHOW TIME
\end{lstlisting}

\subsubsection*{Description}

Displays the current simulation time.

\subsubsection*{Example}

\begin{lstlisting}
SHOW TIME
\end{lstlisting}

Simulation time starts at zero. While the program is running, one real second
advances simulation time by the current multiplier.

\subsection{\texttt{SET TIME MULTIPLIER}}

\subsubsection*{Syntax}

\begin{lstlisting}
SET TIME MULTIPLIER <multiplier>
\end{lstlisting}

\subsubsection*{Description}

Sets the number of simulation-time units added on each automatic one-second
advance. The multiplier must be a non-negative integer. A multiplier of zero
pauses automatic simulation-time advancement.

\subsubsection*{Examples}

\begin{lstlisting}
SET TIME MULTIPLIER 5
SET TIME MULTIPLIER 0
\end{lstlisting}

\subsection{\texttt{FAST FORWARD}}

\subsubsection*{Syntax}

\begin{lstlisting}
FAST FORWARD <delta>
\end{lstlisting}

\subsubsection*{Description}

Immediately increases the current simulation time by the supplied
non-negative delta. The current multiplier is not applied to this operation.
Every intermediate simulation step is processed, allowing agents to act once
per step.

\subsubsection*{Example}

\begin{lstlisting}
FAST FORWARD 30
\end{lstlisting}

\subsection{\texttt{SHOW GRAPH}}

\subsubsection*{Syntax}

\begin{lstlisting}
SHOW GRAPH <symbol>
\end{lstlisting}

\subsubsection*{Description}

Displays the instrument's traded-price and order-book midpoint histories
alongside its exact hidden fundamental-value history. The three series are
shown as separately labelled lines against simulation time. A midpoint is
recorded only while both sides of the order book are available.

\subsubsection*{Example}

\begin{lstlisting}
SHOW GRAPH AAPL
\end{lstlisting}

\subsubsection*{Possible errors}

The command fails if the specified instrument does not exist.

\end{document}
